\section{Introduction}
\label{sec:intro}

% Wu & Tu's introduction is about 1.2 columns and follows a fixed rhythm: (1) the field moved
% away from syntax; (2) transformers won; (3) they were designed heuristically; (4) we derive
% one instead; (5) what we find; (6) what we hope. Match the rhythm -- it is what makes the
% paper read as the sequel it is -- but replace the content with the causal question.

\TODO{draft. Written after the experiments, so the framing matches what was measured.}

\paragraph{Paragraph plan.}
\begin{enumerate}
  \item \citet{wu2023probabilistic} derive a transformer encoder rather than designing one: a
        CRF over dependency structure whose mean-field updates, unrolled, reproduce the
        computation graph of a transformer, with attention emerging as a posterior over
        syntactic heads rather than being posited.
  \item That model is an encoder. Autoregressive generation --- the setting in which
        transformers are actually used --- has no counterpart in the framework, and
        \citet{wu2023probabilistic} name a decoder as an open extension.
  \item Masking the factors does not produce one. \Cref{sec:obstruction} states why: causality,
        a single global energy and contextuality cannot hold together.
  \item Our model: a directed chain of conditional CRFs in which the prefix enters as a frozen
        condition, so the anti-causal term is never formed rather than removed afterwards.
  \item The output mechanism, which is the part to check: the language-model head introduces
        \emph{no} parameter. The word variable, already in the graph, is simply unclamped.
  \item What we find empirically, and the framing --- structure versus weight sharing, not
        versus perplexity.
\end{enumerate}

\paragraph{One principle, two results.}
\NOTE{This paragraph is the one that keeps the paper from reading as two loosely related
contributions. Without it, \Cref{sec:obstruction} and \Cref{sec:output} are an impossibility
theorem and an engineering trick that happen to share a paper.}
Both of our main results are refusals, and they refuse for the same reason: the factor graph is
the model, and it constrains what may be added to it. \Cref{sec:obstruction} says one may not
keep a single global energy and still be causal. \Cref{sec:output} says one may not add an
output matrix that names no factor. Taking the graph seriously in the first case forces the
chain of conditional CRFs; taking it seriously in the second forces tied embeddings, an output
bias equal to the word unary term, and the identity of the encoder and the decoder as two
conditioning patterns of one model.

\paragraph{Contributions.}
\begin{itemize}
  \item A statement of what a causal probabilistic transformer cannot be
        (\Cref{sec:obstruction}), and a construction satisfying the remaining constraints
        (\Cref{sec:chain}).
  \item A graph-faithful output mechanism (\Cref{sec:output}) whose consequences --- forced
        embedding tying, the output bias, one model for both directions --- are derived rather
        than chosen.
  \item An exact readout (\Cref{prop:exact-readout}): the per-position subgraph is a tree, so
        sum-product is exact, and along the sequence it is a single causal
        \texttt{logcumsumexp} scan that preserves the layer-parallel schedule.
  \item An empirical study (\Cref{sec:experiments}) that separates the contribution of
        structure from that of weight sharing by comparing against a Looped transformer as well
        as a standard decoder.
\end{itemize}

\paragraph{What this paper does not claim.}
We do not claim to improve on a standard decoder in perplexity, and the experiments are not
designed to. \citet{wu2023probabilistic} state that the aim of the model is to inform
transformers rather than compete with them, and report degradation beyond roughly $10^5$
sentences; we inherit both the framing and the regime. There is in addition an honest
rank-$\nlab$ bottleneck: all predictive information passes through the label variables. The
question we ask is whether syntactic structure carries information that weight sharing alone
does not, which is a question about the gap to a Looped transformer rather than the gap to a
standard decoder.
